\chapter{Jeux}\label{ch:iv1:games}

\begin{ivtrains}
Réfléchir à ce que fera l'autre camp. Les jeux de retrait (take-away) se ramènent à une table de positions perdantes et aux sommes de Nim, les jeux séquentiels à l'induction à rebours, et les jeux simultanés à la règle selon laquelle une stratégie mixte doit laisser l'adversaire indifférent. Des jeux de poker simplifiés et des enchères réunissent ces outils.
\end{ivtrains}

\begin{question}[$\star$ \iqroles{QR, QT, Trader, Analyste} \iqfirm{toutes} \iqlevel{junior}]\label{iq:iv1:games:1}
Il y a un tas de 97 cailloux. Deux joueurs jouent à tour de rôle en retirant entre 1 et 6 cailloux, et celui qui prend le dernier caillou gagne. Devez-vous jouer en premier, et quel est votre premier coup~? Comment les choses changent-elles si celui qui prend le dernier caillou perd~?
\end{question}

\begin{question}[$\star$ \iqroles{QR, QT, Trader} \iqfirm{toutes} \iqlevel{junior}]\label{iq:iv1:games:2}
Il y a deux tas, de 5 et de 8 cailloux. Deux joueurs jouent à tour de rôle, et chaque coup retire un nombre strictement positif de cailloux d'un seul tas. Celui qui prend le dernier caillou gagne. Qui gagne, et comment~? Qu'en est-il si les tas sont égaux au départ~?
\end{question}

\begin{question}[$\star$ \iqroles{QR, QT, Trader, Analyste} \iqfirm{toutes} \iqlevel{junior}]\label{iq:iv1:games:3}
Vous et moi montrons chacun une pièce au même instant. Si les deux montrent face, je vous paie 3~\$. Si les deux montrent pile, je vous paie 1~\$. S'ils diffèrent, vous me payez 2~\$. Que doit faire chacun de nous, et quelle est la valeur de ce jeu pour vous par manche~?
\end{question}

\begin{question}[$\star$ \iqroles{QR, QT, Trader, Analyste} \iqfirm{toutes} \iqlevel{junior}]\label{iq:iv1:games:4}
Chacun, dans une salle, écrit un nombre entier de 0 à 100. Le gagnant est celui dont le nombre est le plus proche des deux tiers de la moyenne de tous les nombres. Si tout le monde est pleinement rationnel et sait que tous les autres le sont aussi, quel nombre est écrit~? Qu'écririez-vous dans une vraie salle, et pourquoi~?
\end{question}

\begin{question}[$\star\star$ \iqroles{QR, QT, Trader} \iqfirm{toutes} \iqlevel{junior}]\label{iq:iv1:games:5}
Dans un jeu de pierre, papier, ciseaux, une victoire avec la pierre rapporte 1~\$, une victoire avec le papier 2~\$, une victoire avec les ciseaux 3~\$, et une égalité ne rapporte rien. Quelle est la stratégie mixte d'équilibre~? Quelle option est jouée le plus souvent, et pourquoi cela semble-t-il d'abord étrange~?
\end{question}

\begin{question}[$\star\star$ \iqroles{QR, QT, Trader, QD} \iqfirm{toutes} \iqlevel{junior}]\label{iq:iv1:games:6}
Il y a trois tas, de 3, 5 et 7 cailloux. Chaque coup retire un nombre strictement positif de cailloux d'un seul tas, et celui qui prend le dernier caillou gagne. Devez-vous jouer en premier~? Listez tous les premiers coups gagnants, et démontrez la règle que vous utilisez.
\end{question}

\begin{question}[$\star\star$ \iqroles{QR, QT, QD} \iqfirm{toutes} \iqlevel{junior}]\label{iq:iv1:games:7}
Dans un seul tas, un joueur peut retirer 1, 3 ou 4 cailloux à chaque coup, et celui qui prend le dernier caillou gagne. Calculez les valeurs de Grundy pour des tas allant jusqu'à 8 cailloux. Maintenant, le jeu se joue sur deux tas de ce type, de 10 et de 6, et chaque coup est un coup dans un seul tas. Qui gagne, et avec quel coup~?
\end{question}

\begin{question}[$\star\star$ \iqroles{QR, QT, Trader, Analyste} \iqfirm{toutes} \iqlevel{junior}]\label{iq:iv1:games:8}
Cinq pirates, classés de A à E, se partagent 100 pièces d'or. Le pirate vivant le mieux classé propose un partage, et tous les pirates vivants votent. Si au moins la moitié des pirates l'acceptent, le partage est adopté~; sinon, le proposant est jeté par-dessus bord, et le suivant propose. Chacun tient d'abord à sa survie, puis aux pièces, et préfère ensuite voir quelqu'un jeté par-dessus bord. Que propose A~? Comment cela change-t-il si une proposition doit recueillir plus de la moitié des voix~?
\end{question}

\begin{question}[$\star\star$ \iqroles{QR, QT, Trader} \iqfirm{teneur de marché, hedge fund} \iqlevel{junior}]\label{iq:iv1:games:9}
Deux traders se partagent 1~\$ par offres alternées~: A propose un partage, B accepte ou refuse~; si B refuse, B fait une offre, puis A accepte ou refuse, et ainsi de suite. Chaque tour de retard multiplie la valeur de l'argent par $0.9$. S'il y a au plus trois tours, et si après un refus au troisième tour les deux joueurs ne reçoivent rien, que reçoit A~? Et s'il n'y a pas de limite au nombre de tours~?
\end{question}

\begin{question}[$\star\star$ \iqroles{QR, QT, Trader} \iqfirm{toutes} \iqlevel{junior}]\label{iq:iv1:games:10}
Un tireur de penalty vise à gauche ou à droite, et un gardien plonge à gauche ou à droite, au même instant. La probabilité de marquer lorsque le tireur vise à gauche est $0.6$ si le gardien plonge à gauche, et $0.95$ si le gardien plonge à droite. Lorsque le tireur vise à droite, elle vaut $0.9$ si le gardien plonge à gauche, et $0.5$ si le gardien plonge à droite. Déterminez l'équilibre et la probabilité de marquer. Le tireur s'améliore ensuite, et le $0.9$ devient $0.95$. Vise-t-il à droite plus ou moins souvent~?
\end{question}

\begin{question}[$\star\star$ \iqroles{QR, QT, Trader, Analyste} \iqfirm{teneur de marché, crypto} \iqlevel{junior}]\label{iq:iv1:games:11}
$n$ enchérisseurs ont des valeurs privées indépendantes et uniformes sur $[0, 1]$. Dans une enchère au second prix, que doit enchérir chacun, et pourquoi~? Dans une enchère au premier prix, déterminez l'enchère symétrique d'équilibre. Quel revenu le vendeur espère-t-il dans chacune~?
\end{question}

\begin{question}[$\star\star$ \iqroles{QR, QT, Trader, Analyste} \iqfirm{hedge fund, teneur de marché} \iqlevel{junior}]\label{iq:iv1:games:12}
Une tâche vaut 4~\$ à chacune des $n$ personnes, et elle est réalisée si au moins l'une d'elles s'en charge. S'en charger coûte 1~\$ à celle qui le fait. Elles décident toutes au même instant, sans se concerter. Dans l'équilibre symétrique, où chacune s'en charge avec la même probabilité, quelle est cette probabilité pour $n = 2$ et pour $n = 3$~? Quelle est la probabilité que personne ne s'en charge, et que se passe-t-il quand $n$ augmente~?
\end{question}

\begin{question}[$\star\star$ \iqroles{QR, QT} \iqfirm{toutes} \iqlevel{junior}]\label{iq:iv1:games:13}
Une tablette de chocolat est une grille de $m$ par $n$ carrés. Un coup choisit un carré restant et le mange, ainsi que tous les carrés situés au-dessus de lui et à sa droite. Le joueur qui mange le carré en bas à gauche perd. Montrez que le premier joueur possède une stratégie gagnante dès que la tablette comporte plus d'un carré, sans chercher quel est le coup gagnant. Donnez la stratégie explicite pour une tablette carrée.
\end{question}

\begin{question}[$\star\star\star$ \iqroles{QR, QT, Trader} \iqfirm{teneur de marché, hedge fund} \iqlevel{junior, confirmé}]\label{iq:iv1:games:14}
Un équipage de pirates, classés du plus haut au plus bas, se partagent 100 pièces d'or. Le pirate vivant le mieux classé propose un partage, et tous les pirates vivants votent, proposant inclus. Si au moins la moitié des pirates l'acceptent, le partage est adopté~; sinon, le proposant est jeté par-dessus bord, et le suivant dans le classement propose. Chaque pirate tient d'abord à sa survie, puis aux pièces, et préfère ensuite voir quelqu'un jeté par-dessus bord. Avec jusqu'à 200 pirates, le premier proposant survit. Que devient le premier proposant avec 201, 202, 203 et 204 pirates~? Avec 300 pirates, combien sont jetés par-dessus bord~?
\end{question}

\begin{question}[$\star\star\star$ \iqroles{QR, QT, Trader} \iqfirm{teneur de marché, hedge fund} \iqlevel{junior, confirmé}]\label{iq:iv1:games:15}
Un jeu à trois cartes se joue avec un as, un roi et une dame. X et Y versent chacun une mise de 1~\$ et reçoivent une carte. X peut checker, et la carte la plus haute remporte alors le pot de 2~\$, ou miser 1~\$. Face à une mise, Y peut se coucher, et X prend alors le pot, ou suivre, et la carte la plus haute remporte le pot de 4~\$. Déterminez l'équilibre~: à quelle fréquence X bluffe-t-il avec la dame, et à quelle fréquence Y suit-il avec le roi~? Quelle est la valeur du jeu pour X~?
\end{question}

\begin{question}[$\star\star\star$ \iqroles{QR, QT, Trader} \iqfirm{toutes} \iqlevel{junior, confirmé}]\label{iq:iv1:games:16}
Deux enchérisseurs ont des valeurs privées indépendantes, uniformes sur $[0, 1]$, et la valeur que le vendeur attache lui-même au bien est 0. Dans une enchère au second prix avec un prix de réserve $r$, personne ne peut l'emporter en dessous de $r$. Quel prix de réserve maximise l'espérance de revenu du vendeur, et quel est ce revenu~? Le meilleur prix de réserve change-t-il avec davantage d'enchérisseurs~? Que fait un prix de réserve sur le surplus total~?
\end{question}

\begin{question}[$\star\star\star$ \iqroles{QR, QT, QD} \iqfirm{toutes} \iqlevel{junior, confirmé}]\label{iq:iv1:games:17}
Des pièces reposent sur les marches d'un escalier, avec 3, 1, 4, 1 et 5 pièces sur les marches 1 à 5. Un coup prend un nombre strictement positif de pièces d'une marche et les pose sur la marche inférieure~; depuis la marche 1, elles quittent l'escalier. Celui qui effectue le dernier coup gagne. Qui gagne, et quel est le premier coup~? Expliquez pourquoi la moitié des marches peut être ignorée.
\end{question}

\begin{question}[$\star\star\star$ \iqroles{QR, QT, QD} \iqfirm{toutes} \iqlevel{junior, confirmé}]\label{iq:iv1:games:18}
Il y a deux tas de cailloux. Un coup retire un nombre strictement positif de cailloux d'un tas, ou le même nombre strictement positif dans les deux tas. Celui qui prend le dernier caillou gagne. Listez les positions perdantes jusqu'à des tas de 15. La position formée des tas de 9 et de 15 est-elle gagnante pour le joueur qui doit jouer, et celle des tas de 10 et de 15~? Quelle est la forme close des positions perdantes~?
\end{question}

\begin{question}[$\star\star\star\star$ \iqroles{QR, QT} \iqfirm{hedge fund, teneur de marché} \iqlevel{confirmé}]\label{iq:iv1:games:19}
Un jeu de poker utilise trois cartes~: un valet, une dame et un roi. Chaque joueur verse une mise de 1~\$ et reçoit une carte. Le joueur 1 checke ou mise 1~\$. Après un check, le joueur 2 checke, pour aller à l'abattage, ou mise 1~\$, et le joueur 1 se couche ou suit. Après une mise, le joueur 2 se couche ou suit. Déterminez la valeur du jeu pour le joueur 1, et décrivez un ensemble de stratégies d'équilibre. Quelles mains bluffent, et quel joueur a l'avantage~?
\end{question}

\begin{question}[$\star\star\star\star$ \iqroles{QR, QT} \iqfirm{hedge fund, teneur de marché} \iqlevel{confirmé}]\label{iq:iv1:games:20}
X et Y versent chacun une mise de 1~\$ et reçoivent chacun une main privée uniforme sur $[0, 1]$, la main la plus haute étant la meilleure. X peut checker, et la main la plus haute remporte alors le pot de 2~\$, ou miser 1~\$. Face à une mise, Y se couche, et X prend le pot, ou suit, et la main la plus haute remporte le pot de 4~\$. Déterminez l'équilibre, sous la forme suivante~: X mise avec les mains très faibles et les mains très fortes, et Y suit avec une main suffisamment forte. Quelle est la valeur du jeu pour X~?
\end{question}
